
        \documentclass[fleqn]{article}      
        \usepackage{amsmath}
        \usepackage{fontspec}
        \usepackage{kotex} % 한국어 지원  

        \begin{document}      
        ```latex
% Easy Problems
\noindent 문제 1: 다음 행렬의 합을 구하시오.
\[
\begin{pmatrix} 1 & 2 \\ 3 & 4 \end{pmatrix} + \begin{pmatrix} 5 & 6 \\ 7 & 8 \end{pmatrix}
\] 
\\\\\\

\noindent 문제 2: 다음 행렬의 곱을 구하시오.
\[
\begin{pmatrix} 1 & 0 \\ 0 & 1 \end{pmatrix} \times \begin{pmatrix} 2 & 3 \\ 4 & 5 \end{pmatrix}
\]
\\\\\\

\noindent 문제 3: 다음 행렬의 스칼라 곱을 구하시오. 
\[
2 \times \begin{pmatrix} 1 & -1 \\ 3 & 0 \end{pmatrix}
\]
\\\\\\

\noindent 문제 4: 다음 행렬의 전치행렬을 구하시오.
\[
\begin{pmatrix} 1 & 2 & 3 \\ 4 & 5 & 6 \end{pmatrix}
\]
\\\\\\

\noindent 문제 5: 행렬의 결정식을 구하시오.
\[
\begin{vmatrix} 1 & 2 \\ 3 & 4 \end{vmatrix}
\]
\\\\\\

% Medium Problems
\noindent 문제 6: 다음 행렬의 차를 구하시오.
\[
\begin{pmatrix} 5 & -3 \\ 2 & 4 \end{pmatrix} - \begin{pmatrix} 1 & 2 \\ -1 & 0 \end{pmatrix}
\]
\\\\\\

\noindent 문제 7: 다음 두 행렬의 곱을 구하시오.
\[
\begin{pmatrix} 1 & 2 \\ 3 & 4 \end{pmatrix} \times \begin{pmatrix} 5 & 6 \\ 7 & 8 \end{pmatrix}
\]
\\\\\\

\noindent 문제 8: 다음 행렬이 가역행렬인지 확인하시오.
\[
\begin{pmatrix} 2 & 1 \\ 5 & 3 \end{pmatrix}
\]
\\\\\\

\noindent 문제 9: 두 행렬의 전치행렬의 곱을 구하시오.
\[
\begin{pmatrix} 1 & 2 \\ 3 & 4 \end{pmatrix}^T \times \begin{pmatrix} 1 & 0 \\ 0 & 1 \end{pmatrix}^T
\]
\\\\\\

\noindent 문제 10: 다음 행렬의 고유값을 구하시오.
\[
\begin{pmatrix} 1 & 2 \\ 2 & 1 \end{pmatrix}
\]
\\\\\\

% Difficult Problems
\noindent 문제 11: 다음 행렬을 이차 방정식으로 표현하고 고유값을 구하시오.
\[
\begin{pmatrix} 3 & 1 \\ 0 & 2 \end{pmatrix}
\]
\\\\\\

\noindent 문제 12: 두 행렬의 합의 결정식을 구하시오.
\[
\begin{pmatrix} 1 & 2 \\ 3 & 4 \end{pmatrix} + \begin{pmatrix} 5 & 6 \\ 7 & 8 \end{pmatrix}
\]
\\\\\\

\noindent 문제 13: 다음 행렬을 Gauss 소거법을 통해 기약 사다리꼴로 변환하시오.
\[
\begin{pmatrix} 2 & 1 & -1 \\ -3 & -1 & 2 \\ -2 & 1 & 2 \end{pmatrix}
\]
\\\\\\

\noindent 문제 14: 다음 행렬의 역행렬을 구하시오.
\[
\begin{pmatrix} 4 & 7 \\ 2 & 6 \end{pmatrix}
\]
\\\\\\

\noindent 문제 15: 다음 행렬의 쌍대행렬을 구하시오.
\[
\begin{pmatrix} 1 & 0 & 2 \\ 0 & 1 & 3 \end{pmatrix}
\]
\\\\\\

% Answers
1. \(\begin{pmatrix} 6 & 8 \\ 10 & 12 \end{pmatrix}\)\\\\\\
2. \(\begin{pmatrix} 2 & 3 \\ 4 & 5 \end{pmatrix}\)\\\\\\
3. \(\begin{pmatrix} 2 & -2 \\ 6 & 0 \end{pmatrix}\)\\\\\\
4. \(\begin{pmatrix} 1 & 4 \\ 2 & 5 \\ 3 & 6 \end{pmatrix}\)\\\\\\
5. \(-2\)\\\\\\
6. \(\begin{pmatrix} 4 & -5 \\ 3 & 4 \end{pmatrix}\)\\\\\\
7. \(\begin{pmatrix} 19 & 22 \\ 43 & 50 \end{pmatrix}\)\\\\\\
8. \(예, 3\)\\\\\\
9. \(\begin{pmatrix} 1 & 2 \\ 3 & 4 \end{pmatrix}\)\\\\\\
10. \(3, -1\)\\\\\\
11. \(4, 1\)\\\\\\
12. \(0\)\\\\\\
13. \(\begin{pmatrix} 1 & 0 & -0.5 \\ 0 & 1 & 0.5 \\ 0 & 0 & 0 \end{pmatrix}\)\\\\\\
14. \(\begin{pmatrix} 3/10 & -7/10 \\ -1/5 & 2/5 \end{pmatrix}\)\\\\\\
15. \(\begin{pmatrix} 1 & 0 \\ 0 & 1 \\ 2 & 3 \end{pmatrix}\)\\\\\\
```      
        \end{document} 
    